% Condensed review; existing citation coverage retained without new source claims.
\section{Literature Review}\label{sec:literature}\label{sec:related}

Salsa20 and ChaCha combine ARX rounds with modular feed-forward to produce 512-bit blocks~\citep{bernstein2008salsa,estream2008,nir2018}; ChaCha's quarter-round was designed to improve diffusion per round~\citep{bernstein2008chacha}.

Chosen-IV methods establish the fixed-key, related-input setting~\citep{englund2007,fischer2008,maitra2016}. Differential cryptanalysis~\citep{biham1991} and probabilistic neutral bits (PNBs)~\citep{aumasson2008} underpin multi-bit and differential-linear attacks on reduced Salsa and ChaCha~\citep{choudhuri2016,dey2017,beierle2020,coutinho2023}. Later PNB studies refine input differences, bias estimates, and neutral-bit selection~\citep{ghafoori2023,dey2024salsa,dey2024tit}. Higher-order analysis~\citep{ghafoori2024hodl}, two-round key--IV criteria~\citep{dey2025tosc}, bit puncturing~\citep{florez2025}, and automated correlation search~\citep{liu2026} extend this line. Multiple-output-difference PNB results are also reported in a preprint~\citep{paul2025}.

Avalanche and bit-independence tests~\citep{webster1986,hernandez2005,madarro2024}, together with NIST's Fourier test~\citep{nist2010}, provide established statistical baselines. Spectral entropy has long been used in time-series analysis~\citep{powell1979,inouye1991}, with later normalized and complexity--entropy representations~\citep{zaccarelli2013,dai2018,mao2018}. A recent cryptographic application reports that unconditional ChaCha20 keystream meets its spectral-entropy criterion~\citep{kasyanov2026}. Here spectral entropy scores ordered HW trajectories of first-order nonce differences; mixed-derivative HW profiles separately probe finalization. A passing unconditional test does not characterize those responses.

Our contribution connects the conditional word-7 rank and mean-singleton identities, paired add/none/xor responses, and frozen-model spectral reconstruction on new keys. Two-round Salsa20 criteria~\citep{dey2025tosc} and higher-order ChaCha analysis~\citep{ghafoori2024hodl} already place low-round structure in cryptanalysis; we claim no novelty for weak two-round diffusion or higher-order differences themselves. Here local identities predict zeros and interactions, interventions isolate finalization effects, and fresh keys test empirical reconstruction. This does not derive the entire entropy gap from one word or improve attack complexity.
